\documentclass{article}
\usepackage{graphicx}
\usepackage{amsmath}
\usepackage{booktabs}
\usepackage{array}
\usepackage{enumitem}
\usepackage{xcolor}
\usepackage{titlesec}
\usepackage{fancyhdr}
\usepackage{listings}

\title{Peer-to-Peer Networks}
\author{0xSaad}
\date{\today}

\setlist[itemize]{topsep=3pt,itemsep=1pt,parsep=0pt}
\setlist[enumerate]{topsep=3pt,itemsep=1pt,parsep=0pt}

\begin{document}

\maketitle

\section{Client/Server and Peer-to-Peer}

\subsection{Client/Server Architecture}

A well-known, powerful, reliable server is a data source. Clients request data from the server.

Examples: WWW (HTTP), FTP, and Web services.

\subsubsection{Client/Server Limitations}
\begin{itemize}
    \item Scalability is hard to achieve.
    \item Presents a single point of failure.
    \item Requires administration.
\end{itemize}

P2P systems try to address these limitations.

\subsection{P2P Computing}

P2P computing is the sharing of computer resources and services by direct exchange between systems.

\begin{itemize}
    \item Resources include information, processing cycles, cache storage, and disk storage for files.
    \item P2P takes advantage of existing computing power, storage, and networking connectivity.
\end{itemize}

\subsection{P2P Architecture}
\begin{itemize}
    \item All nodes are both clients and servers: they provide and consume data.
    \item Any node can initiate a connection.
    \item No centralized data source.
\end{itemize}

\subsubsection{P2P Network Characteristics}
\begin{itemize}
    \item Clients are also servers and routers; nodes contribute content, storage, memory, and CPU.
    \item Nodes are autonomous (no administrative authority).
    \item The network is dynamic: nodes frequently enter and leave.
    \item Nodes collaborate directly, not through well-known servers.
    \item Nodes have widely varying capabilities.
\end{itemize}

\subsubsection{P2P Benefits}
\begin{itemize}
    \item Efficient use of resources: unused bandwidth, storage, and processing power at the edge of the network.
    \item Scalability: consumers also donate resources; aggregate resources grow with utilization.
    \item Reliability: replicas, geographic distribution, and no single point of failure.
    \item Ease of administration: self-organizing nodes, fault tolerance, replication, and load balancing.
\end{itemize}

\subsubsection{P2P Applications}
\begin{itemize}
    \item File sharing: Napster, Gnutella, Kazaa.
    \item Multiplayer games: Unreal Tournament, DOOM.
    \item Collaborative applications: ICQ, shared whiteboard.
    \item Distributed computation: SETI@home.
    \item Ad-hoc networks.
\end{itemize}

Large-scale file sharing may raise issues of copyright infringement.

\section{P2P Search Architectures}

\subsection{Approaches}
\begin{itemize}
    \item Centralized: Napster.
    \item Flooding or Decentralized: Gnutella.
    \item Hybrid (Centralized and Decentralized) with Super-Peers: KaZaA.
    \item Document Routing (Structured): Chord and CAN.
\end{itemize}

\subsection{Napster (Centralized)}

A system for sharing music files using centralized searching and direct peer-to-peer file transfer.

\begin{enumerate}
    \item Users upload their file lists to the Napster server.
    \item Users send keyword queries to the server (artist, song, album, bitrate, etc.).
    \item The server replies with IP addresses of users with matching files.
    \item Users connect directly to another peer to download the file.
\end{enumerate}

\subsubsection{Napster Properties}
\begin{itemize}
    \item A central server can ensure correct search results.
    \item The server is a bottleneck for scalability.
    \item Single point of failure.
    \item Susceptible to denial of service, malicious users, lawsuits, and legislation.
    \item Search is centralized; file transfer is direct (peer-to-peer).
\end{itemize}

\subsection{Gnutella (Flooding)}

A decentralized P2P system that shares any type of files, not only music.

\begin{enumerate}
    \item A peer asks its neighbors for files of interest.
    \item Neighbors forward the query to their neighbors, and so on.
    \item A TTL (Time To Live) field stops messages after a number of hops.
    \item Peers with matching files reply to the requester.
\end{enumerate}

\subsubsection{Gnutella Properties}
\begin{itemize}
    \item Decentralized, with no single point of failure.
    \item Not as susceptible to denial of service.
    \item Cannot ensure correct results.
    \item Flooding makes search distributed but still not scalable.
\end{itemize}

\subsection{KaZaA (FastTrack Network)}

A hybrid of centralized Napster and decentralized Gnutella.

\begin{itemize}
    \item Super-Peers act as local search hubs, like Napster servers for small portions of the network.
    \item Super-Peers are automatically selected by capacity (storage, bandwidth) and availability (connection time).
    \item Users upload file lists to a Super-Peer.
    \item Super-Peers periodically exchange file lists.
    \item Users send file queries to a Super-Peer.
\end{itemize}

\subsection{Free Riding}

File-sharing networks rely on users sharing data.

Two types of free riding:
\begin{itemize}
    \item Downloading but not sharing any data.
    \item Not sharing any interesting data.
\end{itemize}

On Gnutella, the study cited in the slides reports:
\begin{itemize}
    \item 15\% of users contribute 94\% of content.
    \item 63\% of users never responded to a query.
\end{itemize}

\subsection{Anonymity and Freenet}

Napster, Gnutella, and Kazaa do not provide anonymity: users know who they download from, and others know who sent a query.

Freenet is designed to provide anonymity.

\begin{itemize}
    \item Data flows in the reverse path of the query.
    \item It is impossible to know if a user is initiating or forwarding a query, or consuming or forwarding data, in the model described in the slides.
    \item ``Smart'' queries route requests toward the correct peer by incremental discovery.
\end{itemize}

\subsubsection{Freenet Routing Tables}
\begin{itemize}
    \item \texttt{id}: File identifier (e.g., hash of file).
    \item \texttt{next\_hop}: Another node that stores the file ID.
    \item \texttt{file}: File identified by the ID, stored at the local node.
\end{itemize}

Forwarding a query for a file ID:
\begin{enumerate}
    \item If the file is stored locally, stop and forward data back to the upstream requester.
    \item Otherwise, find the closest ID in the routing table and forward to its \texttt{next\_hop}.
    \item If data is not found, report failure; the requester tries the next closest match.
\end{enumerate}

\section{Structured P2P and Distributed Hash Tables}

\subsection{Structured P2P}

Second-generation P2P overlay networks based on a Distributed Hash Table (DHT) interface.

\begin{itemize}
    \item Self-organizing.
    \item Load balanced.
    \item Fault-tolerant.
    \item Scalable guarantees on the number of hops to answer a query, unlike unstructured P2P systems.
\end{itemize}

\subsection{Distributed Hash Tables (DHT)}

A distributed version of a hash table data structure that stores \texttt{(key, value)} pairs.

\begin{itemize}
    \item The key is like a filename; the value can be file contents.
    \item Goal: Efficiently insert, lookup, and delete \texttt{(key, value)} pairs.
    \item Each peer stores a subset of the key-value pairs.
    \item Core operation: Find the node responsible for a key.
    \item Map a key to a node and efficiently route requests to that node.
\end{itemize}

\subsubsection{Directed Searches}
\begin{itemize}
    \item Assign particular nodes to hold particular content, or know where it is.
    \item A request goes toward the node responsible for the content.
    \item Challenges: Evenly distributing responsibilities and adapting to nodes joining, leaving, or failing.
    \item Give responsibilities to joining nodes; redistribute responsibilities from leaving nodes.
\end{itemize}

\subsubsection{Hash Tables and DHTs}
\begin{itemize}
    \item A hash table associates data with keys and hashes a key to find a bucket.
    \item Each bucket is expected to hold \(\text{items}/\text{buckets}\) items.
    \item In a DHT, nodes are the hash buckets: a key is hashed to find its responsible peer.
    \item Data and load are balanced across nodes.
\end{itemize}

\subsubsection{DHT Applications}
\begin{itemize}
    \item File sharing and archival storage.
    \item Databases.
    \item Naming and service discovery.
    \item Chat service.
    \item Rendezvous-based communication.
    \item Publish/Subscribe.
\end{itemize}

\subsubsection{DHT Desirable Properties}
\begin{itemize}
    \item Keys are mapped evenly to nodes (e.g., 10 keys and 5 peers: 2 keys per peer).
    \item Each node maintains information about only a few other nodes.
    \item Messages are routed efficiently.
    \item Node arrivals and departures affect only a few nodes.
\end{itemize}

\subsubsection{DHT Routing Protocols}

DHT is a generic interface with several implementations:
\begin{itemize}
    \item Chord (MIT).
    \item Pastry (Microsoft Research UK, Rice University).
    \item Tapestry and CAN (UC Berkeley).
    \item SkipNet, Kademlia, Viceroy, P-Grid, and Freenet.
\end{itemize}

These systems are called P2P routing substrates or P2P overlay networks.

\section{Chord}

\subsection{Chord API}
\begin{itemize}
    \item Node ID: Unique \(m\)-bit identifier (hash of IP address or another unique ID).
    \item Key: \(m\)-bit identifier (hash of a sequence of bytes).
    \item Value: Sequence of bytes.
\end{itemize}

Operations:
\begin{itemize}
    \item \texttt{insert(key, value)}: Store key/value at \(r\) nodes.
    \item \texttt{lookup(key)}.
    \item \texttt{update(key, newval)}.
    \item \texttt{join(n)}.
    \item \texttt{leave()}.
\end{itemize}

\subsection{Chord Properties}

In a system with \(N\) nodes and \(K\) keys, with high probability, the slides give:
\begin{itemize}
    \item Key responsibility on the order of \(K/N\) per node.
    \item Each node maintains information about \(O(\log N)\) other nodes.
    \item Lookups are resolved in \(O(\log N)\) hops.
    \item No delivery guarantees.
    \item No consistency among replicas.
\end{itemize}

\subsection{Chord Identifier Ring}

Chord uses a circular \(m\)-bit ID space for both keys and nodes:
\[
0,1,\ldots,2^m-1.
\]

\begin{itemize}
    \item Node ID: SHA-1 of an IP address (as shown in the slides).
    \item Key ID: SHA-1 of a key.
    \item A key maps to the first node whose ID is equal to or follows the key ID on the ring.
    \item Each node is responsible for \(O(K/N)\) keys; \(O(K/N)\) keys move when a node joins or leaves.
\end{itemize}

Example: \(m=6\), so \(2^6=64\) possible IDs (0--63).

With nodes \(\{1,8,14,21,32,38,42,48,51,56\}\), the mappings illustrated in the slides include:
\[
K10\to N14,\quad K24\to N32,\quad K30\to N32,\quad
K38\to N38,\quad K54\to N56.
\]

\subsection{Chord State and Lookup}

\subsubsection{Basic Chord}
\begin{itemize}
    \item Each node knows its Successor and Predecessor (for ring management).
    \item Lookup requests are forwarded around the ring through Successor pointers.
    \item Requires \(O(N)\) hops.
\end{itemize}

\subsubsection{Finger Table}
\begin{itemize}
    \item Each node knows \(m\) other positions on the ring.
    \item Finger \(i\) of node \(n\) points to the Successor of \(n+2^i\).
    \item Lookup follows the closest preceding fingers and then the Successor.
    \item Requires \(O(\log N)\) hops.
\end{itemize}

\[
\operatorname{finger}_i(n)=\operatorname{successor}\bigl((n+2^i)\bmod 2^m\bigr),
\qquad i=0,\ldots,m-1.
\]

Finger table for \(N8\) when \(m=6\), as in the slides:

\begin{center}
\begin{tabular}{ccc}
\toprule
\(i\) & \(N8+2^i\) & Successor \\
\midrule
0 & 9  & N14 \\
1 & 10 & N14 \\
2 & 12 & N14 \\
3 & 16 & N21 \\
4 & 24 & N32 \\
5 & 40 & N42 \\
\bottomrule
\end{tabular}
\end{center}

\subsection{Chord Join and Routing}

In the identifier space \([0,7]\), \(m=3\).

\begin{itemize}
    \item Node \(n1\) joins first: every Successor-table entry points to \(n1\).
    \item Node \(n2\) joins: the Successor tables change.
    \item Nodes \(n0\) and \(n6\) join: active nodes are \(\{n0,n1,n2,n6\}\).
    \item Items \(f1\) and \(f7\): \(f1\) maps to \(n1\), and \(f7\) maps to \(n0\).
\end{itemize}

Example Successor tables for the four active nodes (each row lists entries for \(i=0,1,2\)):

\begin{center}
\begin{tabular}{cl}
\toprule
Node & Successor-table entries \\
\midrule
\(n0\) & \(n1,\ n2,\ n6\) \\
\(n1\) & \(n2,\ n6,\ n6\) \\
\(n2\) & \(n6,\ n6,\ n6\) \\
\(n6\) & \(n0,\ n0,\ n2\) \\
\bottomrule
\end{tabular}
\end{center}

Upon receiving a query for an item ID:
\begin{enumerate}
    \item Check if the item is stored locally.
    \item Otherwise, forward to the largest preceding node in the Successor table.
    \item Continue until the responsible node is reached.
\end{enumerate}

Example from the slides: Query for item \(7\) from \(n1\) is routed toward \(n6\), then \(n0\).

\subsection{Network Locality}

Nodes close on the Chord ring can be far apart in the physical network.

\section{Content Addressable Network (CAN)}

\subsection{Document Routing -- CAN}

CAN associates each node and item with a unique identifier in a \(d\)-dimensional space.

Goals:
\begin{itemize}
    \item Scale to hundreds of thousands of nodes.
    \item Handle rapid arrival and failure of nodes.
\end{itemize}

Properties:
\begin{itemize}
    \item Routing table size: \(O(d)\).
    \item File location in at most \(d\,n^{1/d}\) routing steps, where \(n\) is the total number of nodes, as stated in the slides.
\end{itemize}

\subsection{CAN Two-Dimensional Space}
\begin{itemize}
    \item The space is divided between nodes; all nodes together cover the entire space.
    \item Each node covers a square or rectangular area (ratios 1:2 or 2:1 in the slide example).
    \item The first node, \(n1:(1,2)\), covers the entire space.
    \item When \(n2:(4,2)\) joins, the space is divided between \(n1\) and \(n2\).
    \item Later nodes join: \(n3:(3,5)\), \(n4:(5,5)\), and \(n5:(6,6)\).
\end{itemize}

\subsubsection{Data Placement}

Each item is stored by the node that owns its mapped area in the coordinate space.

\begin{center}
\begin{tabular}{ccc}
\toprule
Item & Coordinates & Owning Node \\
\midrule
\(f1\) & \((2,3)\) & \(n1\) \\
\(f2\) & \((5,1)\) & \(n2\) \\
\(f3\) & \((2,1)\) & \(n1\) \\
\(f4\) & \((7,5)\) & \(n5\) \\
\bottomrule
\end{tabular}
\end{center}

\subsection{CAN Query Routing}
\begin{itemize}
    \item Each node knows its neighbors in the \(d\)-dimensional space.
    \item Forward a query to the neighbor closest to the query ID.
    \item Slide example: \(n1\) queries item \(f4\).
    \item CAN can route around some failures; some failures require local flooding.
\end{itemize}

\section{P2P Web}

Works like an ordinary Web server, with these differences:
\begin{itemize}
    \item Distributed storage.
    \item Automatic data replication.
    \item Load balancing of users.
    \item When the load is too high, add another server and store another copy.
\end{itemize}

\end{document}
